\chapter{绪论}

\section{背景和意义}

舵机是一类典型的机电执行机构，常用于将控制系统输出的指令转化为可测量的角位移、角速度或力矩响应。在航空航天、自动控制和工业装备等场景中，舵机通常处在控制链路的末端，其动态响应、稳态误差、重复性和通信稳定性会直接影响上层系统的控制效果。对于精密舵机而言，单次动作完成只是基本要求；在连续指令、不同激励形式和不同测试流程下，系统还需要保持可观测、可分析和可追溯的工作状态。国内伺服测试平台和工业测试系统研究表明，执行机构测试既涉及数据采集，也涉及测试对象、控制链路、数据处理和结果管理之间的流程组织\cite{li2014servotest,xu2021industry4test,liu2022industrialsoftware}。

在实际测试过程中，舵机测试软件需要同时处理设备连接、指令发送、反馈解析、曲线显示、性能计算和结果保存等任务。这些任务对应的处理频率存在差异。串口反馈、帧解析和曲线刷新属于高频数据流，要求软件能够持续接收并快速处理；参数输入、测试项选择、结果查看和数据导出属于低频控制流，要求界面保持清晰和可操作。将这些任务简单集中在同一执行上下文中，容易出现界面阻塞、数据堆积、响应滞后和测试状态不一致等问题。数据采集软件和在线分析系统的相关研究也指出，采集、处理、展示和分析之间需要明确边界，软件架构本身会影响实时数据流的稳定性和可维护性\cite{stansbury2023autodidaqt,turkmanovic2023highspeed}。

本课题面向精密舵机测试场景，设计并实现一套桌面测试软件。该软件运行于 Windows 环境，采用 RS422 串口链路与舵机控制器通信，前端使用 Vue3 构建参数输入和数据展示界面，后端基于 Electron 主进程、预加载脚本和 Worker 机制组织串口通信、数据处理和测试执行逻辑。所设计的软件系统以分层架构为基础，将协议解析、实时显示、通信性能测试、控制性能预测试、时域性能测试、频域性能测试和结果保存组织到同一测试流程中。

从工程实现角度看，本课题的意义主要体现在三个方面。第一，针对高频反馈数据和低频界面交互并存的问题，采用分层架构和 Worker 机制划分职责，降低界面层与数据处理层之间的耦合。第二，针对舵机通信协议存在差异的问题，将控制帧生成、反馈帧解析、校验和结构化数据表达作为独立通信层设计，避免测试逻辑直接依赖原始字节流。第三，针对性能测试参数来源不清和测试结果难以复查的问题，将预测试、统计参数推导、时域及频域指标提取和数据保存纳入统一流程，使测试过程能够被解释、记录和复核。

论文写作口径以当前软件完成状态为准。本系统定位为精密舵机桌面测试软件，早期需求设想中的其他测试内容只作为需求演化背景。当前正文重点围绕四类已经落地的测试功能展开：通信性能测试、控制性能预测试、时域性能测试和频域性能测试。长时间稳定性、完整端到端延迟等缺少独立实测支撑的指标，仅在设计目标、实现约束或后续工作中讨论，相关结论不列为已验证成果。

\section{国内外研究现状}

本课题涉及执行机构测试、工业数据采集软件、串口通信可靠性和自动测试方法等多个方向。直接以“基于 Electron 的多协议舵机测试软件”为主题的文献较少，因此本节围绕本文要解决的问题，将现有研究归纳为四个方面，并说明各类研究对本文的启发和不足。

\subsection{执行机构测试平台与硬件在环研究}

执行机构测试平台研究主要关注如何在可控条件下复现执行机构的工作过程，并通过测试台或硬件在环系统测量对象的动态特性、负载响应和控制性能。李敬孔围绕伺服系统测试平台的设计与实现，讨论了测试对象分析、平台功能组织和结果处理等内容，为本文理解伺服类对象的测试软件边界提供了国内研究参考\cite{li2014servotest}。Anastasopoulos 和 Hornung 针对无人机舵机执行器设计实时测试台，关注执行器在测试环境下的动态响应测量\cite{anastasopoulos2018uavbench}。Dell'Amico 面向航空执行机构提出硬件在环空气载荷模拟系统，通过更接近真实工况的载荷复现来验证执行机构性能\cite{dellamico2024airloadhil}。

国内围绕单对象控制实验平台也积累了较多工作。林立等设计了集成 DSP 控制器和上位机监控界面的异步电机直接转矩控制实验平台，周晓华等利用 Matlab/Simulink 构建永磁同步电机矢量控制仿真实验，张惠国等则将低压直流无刷电机控制系统组织为综合性实验平台\cite{lin2025dtcplatform,zhou2020pmsmexperiment,zhang2016bldcexperiment}。这些研究虽然主要服务于教学验证或单对象实验，但都强调了控制回路、反馈采集、上位机界面和实验流程之间的一体化组织。

这些研究表明，执行机构测试已经从单次测量逐渐发展为包含工况设置、激励生成、反馈采集和性能分析的系统工程。测试平台的价值体现在数据获取、过程稳定性、重复性和可解释性等方面。对本文而言，这一认识具有直接启发意义：舵机测试软件应承担测试平台中的软件组织层作用。

在执行机构控制策略方向，汪春阳针对电液位置伺服系统研究了滑模、有限时间和预设性能控制，并通过实验验证控制效果；王超则面向重载电液伺服驱动机器人研究末端拖拽控制与阻抗控制策略\cite{wang2025electrohydraulicsmc,wang2025hydraulicdrag}。这类工作说明，国内研究已经在特定执行机构的控制品质提升和实验验证方面形成较深积累，但其重点仍然集中在单对象控制器设计，而不是多协议桌面测试软件的流程组织。

不过，已有执行机构测试平台研究也存在与本文场景不完全匹配之处。许多研究重点放在专用测试台、硬件加载装置或特定执行器对象上，软件部分往往服务于固定实验系统。对于桌面测试软件内部如何组织多协议解析、高频数据流、测试项切换和结果保存，相关讨论相对有限。本文的切入点正是在执行机构测试平台思想的基础上，进一步讨论面向精密舵机测试的软件架构和测试流程组织问题。

\subsection{工业桌面测试系统架构与实时数据采集研究}

设备测试软件需要处理持续到达的数据流，也需要响应用户操作、展示曲线并保存结果。随着测试对象和数据规模增加，传统单线程或单流程界面程序很难同时保证实时性和可维护性。徐景辉等从工业 4.0 条件下的测试系统出发，指出测试系统正在向自动化、网络化和信息融合方向发展\cite{xu2021industry4test}。刘务等关于工业软件测试验证体系的研究强调，工业软件需要通过规范化验证链路保证可靠性\cite{liu2022industrialsoftware}。Stansbury 和 Lanzara 提出的 AutodiDAQt 将数据采集与在线分析结合，强调科学数据采集软件应支持并发、可复现、采集过程回溯和自动化界面生成\cite{stansbury2023autodidaqt}。Turkmanovi\'c 等针对远程高速数据采集提出高性能软件架构，指出低延迟传输和实时显示不仅依赖硬件链路，也依赖软件端任务划分与数据处理方式\cite{turkmanovic2023highspeed}。

国内研究也为桌面测试软件的技术选型提供了参考。张子磐围绕 Qt 测控软件开发讨论了桌面测控系统中的采集、控制和界面集成问题\cite{zhang2017qtflowmeter}。高雨等讨论了基于 Electron 框架构建跨平台数据采集桌面应用的方法，说明 Web 技术栈在桌面数据采集软件中具有界面开发效率和跨平台部署优势\cite{gao2025electron}。在此基础上，前端界面采用组件化方式组织不同测试面板和结果区域，使通信选择、参数输入、控制与数据显示能够围绕同一测试流程协同展开。这些研究虽不直接面向舵机测试，但能够支撑本课题所设计软件在桌面框架选择与界面组织方式上的合理性。

总体来看，现有研究已经形成若干共识：采集、处理、显示和保存应尽量分层；高频数据处理不宜直接压在界面线程；测试数据需要以结构化形式保存，便于后续追溯和分析。所设计的软件系统吸收这些思路，将渲染进程、主进程、预加载桥接和 Worker 模块结合起来，面向舵机测试过程组织高频反馈数据和低频测试控制逻辑。

\subsection{RS422/UART 串口通信与可靠传输研究}

通信链路是舵机测试软件的基础。只有控制指令能够正确发送、反馈数据能够稳定接收、帧边界能够识别、字段能够解析并完成校验，后续测试算法和可视化结果才有意义。所设计的软件系统通过 RS422 与舵机控制器连接，论文重点关注串口链路之上的协议解析、帧校验和结构化数据表达。

在公开研究中，直接讨论 RS422 多协议舵机测试软件的文献较少。与本文更接近的中文研究多分布在测控软件和串口通信实现方向。张子磐关于 Qt 测控软件开发的研究说明，桌面测控系统通常需要在设备通信、数据采集、界面显示和结果管理之间建立清晰接口\cite{zhang2017qtflowmeter}。吴浩松等研究了基于数据分包的嵌入式串口通信方法，指出在连续数据流中合理的分包、缓存和校验机制有助于提高通信稳定性\cite{wu2025serialpacket}。

在可靠性方面，串口链路的稳定性不仅取决于物理接口，也取决于上层软件对帧边界、缓存区、校验结果和异常数据的处理方式。对于舵机测试软件而言，通信模块需要把原始字节流转换为结构化反馈数据，并在解析失败、半包、粘包或校验错误时给出明确处理结果。上述研究为所设计的软件系统建立 RS422 反馈帧解析、缓存重组和错误上报机制提供了基础参考。

这些研究对本文的启发在于，RS422 可以被视为底层传输链路，测试软件需要在上层建立可维护的协议处理机制。所设计的软件系统围绕同步字节、帧长度、控制量字段、反馈量字段和校验规则组织协议解析流程，并将解析后的反馈数据转化为统一结构，供可视化和测试算法复用。现有研究更多集中在串口基础机制、硬件实现或单一协议场景，对桌面测试软件中多协议配置、半包及粘包处理、解析错误上报和业务层数据复用讨论不足，这也是本文需要进一步展开的内容。

\subsection{自动测试、时域/频域评估与实验设计研究}

舵机测试软件除了要解决通信和界面问题，还要回答“如何测试”的问题。若测试参数全部依靠人工经验设定，测试结果容易受到人为差异影响；若测试流程只追求覆盖所有点位和频率，又会导致耗时过长。因此，自动测试方法需要在测试质量、统计可靠性和执行时间之间取得平衡。

在动态特性评估方面，Cajetinac 等通过 PWM 控制和频率特性辨识研究气动执行器，说明执行机构频域测试需要关注激励信号、采样过程和结果计算之间的关系\cite{cajetinac2012pwm}。此外，测试流程还需要在覆盖范围、重复次数和总测试时间之间取得平衡。国内工业测试系统和工业软件测试验证研究也指出，测试系统应兼顾测试流程规范、结果可追溯和工程现场使用效率\cite{xu2021industry4test,liu2022industrialsoftware}。

这些研究共同说明，测试软件需要把测试参数、激励设计、数据采集和指标计算组织为可解释流程。所设计的软件系统在时域测试和频域测试中采用预测试作为参数来源：先估计噪声水平、响应时间等中间量，再用于正式测试中的窗口长度、重复次数、采样时长或频点配置。频域测试中，系统通过正弦激励和相关分析提取幅值与相位；时域测试中，系统围绕调节时间、超调量、稳态误差和对称性等指标组织结果。本文的重点落在这些方法进入桌面测试软件后的执行流程、实时显示和结果保存方式。

综上，现有研究已经在执行机构测试平台、实时数据采集软件、串口可靠通信和自动测试方法方面形成了较多成果，但这些成果分散在不同问题域中。对于本文面向的精密舵机测试软件，仍需要将协议解析、分层架构、高频数据处理、四类测试功能和结果闭环统一起来。本课题正是在这一交叉位置展开软件系统的设计与实现。

\section{本文主要工作}

围绕精密舵机测试软件的设计与实现，本文从系统结构、通信链路、测试流程和结果管理几个方面展开工作。系统采用 Electron、Vue3 和 TypeScript 技术栈构建桌面应用，前端负责参数输入、测试项选择、数据显示和结果交互，后端通过主进程、预加载脚本和 Worker 模块组织串口通信、测试任务和渲染数据整理。这样的结构将界面交互、高频数据处理和测试逻辑分离，使软件在处理连续反馈数据时仍能保持清晰的模块边界。

在通信与测试功能实现方面，所设计的软件系统建立了 RS422 串口通信链路和协议解析机制。软件能够进行串口端口选择、连接管理、控制帧发送、反馈帧接收、字段解析和校验处理，并将原始反馈转化为上层模块可复用的结构化数据。在此基础上，系统围绕当前已经落地的四类测试功能组织核心流程，包括通信性能测试、控制性能预测试、时域性能测试和频域性能测试。通信性能测试用于统计发送和接收过程中的成功率、采样间隔和置信区间；控制性能预测试用于估计噪声水平、响应时间和运动速度等参数；时域性能测试用于分析调节时间、超调量和稳态误差等动态指标；频域性能测试则通过扫频和相关法估计幅频、相频特性，并形成伯德图等结果展示。

为了使测试过程便于观察和复查，所设计的软件系统还提供实时可视化和结果管理功能。系统提供四通道遥测数据显示、示波器曲线、时域图表和频域伯德图，并支持测试结果保存和历史数据加载。围绕该软件系统，工程中设置了串口解析、协议合规、信号发生器、通信性能测试、时域测试和频域测试等单元测试，也包含面向真实硬件的后端冒烟测试和实机测试脚本。本文的主要创新点集中在软件架构、通信链路和测试流程组织，上述测试用于支撑系统实现的正确性和可维护性。

\section{论文组织结构}

绪论部分从精密舵机测试的软件化需求出发，说明课题背景、研究意义和相关研究基础。该部分在研究现状梳理的基础上，将执行机构测试平台、工业桌面测试系统、串口通信和自动测试方法等分散问题归纳到本文的软件设计任务中，由此明确后续需求分析需要回答的对象、边界和目标。

需求分析与相关技术部分将测试对象、使用场景和最终软件状态转化为可设计的约束条件。通过梳理串口通信、参数配置、实时显示、四类测试和结果管理等功能要求，可以避免总体设计停留在技术栈介绍层面，使后续架构划分能够直接服务于实际测试流程。同时，相关技术的讨论也为 Electron、Vue3、TypeScript、Worker 和串口通信机制的选择提供依据。

系统总体设计部分在需求约束的基础上确定软件的分层结构和数据流组织方式。该部分重点说明渲染进程、主进程、预加载脚本和 Worker 模块如何分担人机交互、任务调度、高频数据处理和测试计算等职责，其意义在于先建立清晰的模块边界，再为通信协议、基础控制和测试算法的展开提供统一框架。

通信协议与基础控制模块设计部分沿着总体架构中的数据通路继续展开，重点处理软件与舵机控制器之间的数据交换问题。RS422 通信链路、协议抽象、控制帧生成、反馈帧解析、使能检测和信号发生器等内容共同构成测试流程的底层支撑。只有先把硬件反馈转化为结构化数据，并提供可复用的控制与激励能力，后续测试功能才能稳定获得输入并执行指标计算。

测试功能与算法设计部分在通信和基础控制能力之上组织具体测试流程。通信性能测试用于先检查数据链路质量，控制性能预测试用于估计噪声水平、响应时间和运动速度等中间参数，时域性能测试和频域性能测试再利用这些参数确定激励设置、采样窗口和指标计算方式。这样的安排体现了软件测试流程由链路确认、参数估计到正式性能分析的推进关系。

系统实现与测试验证部分将前述需求、架构、通信基础和测试流程落实到工程代码中，并通过单元测试、集成测试和实机测试检查软件是否能够贯通运行。该部分的意义在于把设计说明转化为可验证的软件事实，既说明关键模块如何实现，也为全文总结提供实现效果和问题边界方面的依据。

总结与展望部分在实现和验证结果的基础上回顾全文工作，归纳精密舵机测试软件在分层架构、串口通信、测试流程组织和结果管理方面完成的内容。同时，该部分结合前文验证中尚未充分覆盖的问题，提出后续在长时间稳定性验证、更多协议适配和测试报告自动化方面的改进方向，使论文从现有实现自然过渡到后续完善工作。
